% October 13, 2026 build (agent draft): abstract at the 7.9 production point.
% Numbers: 7p9167_v6/table.md (from-2023 GE transitions, each policy against
% its own no-policy path), risk0_baselines (1.641 -> 1.905).
\begin{abstract}
\noindent
Housing costs are often blamed for low fertility, and the remedies proposed
range from easier mortgage credit to building more homes. This paper studies
which of them can raise fertility in a life-cycle model in which households
choose when to have children, how much space to occupy, and whether to rent
or own it. Children need space, the largest homes must be bought with a down
payment, and earnings are subject to uninsurable risk. I calibrate the model
to the United States in 2007, attribute the subsequent fall in fertility to
changes in the preference for children, and treat 2023 as a point on a
transition to a smaller population. From the 2023 economy, a 10 percent
expansion of the housing supply raises the completed fertility of the cohort
entering adulthood by 4.8 percent, and the effect persists for later cohorts.
Removing the down payment raises births by 1.9 percent on impact but leaves
completed fertility unchanged: it moves first births earlier and buys
ownership of small homes. The same credit restricted to parents raises
completed fertility by 6.3 percent. Earnings risk matters as much as housing:
a mean-preserving reduction in its standard deviation from 0.48 to 0.35 raises
completed fertility by 16 percent. A child commits a household to paying for
family space every period, and policies raise fertility when they lower that
cost or make it easier to bear.
\end{abstract}
